% !TeX root = ../Western Philosophy.tex
\section{近代哲学的背景}

\subsection{近代历史的奠基}
Discart（1596）16世纪末期出生，一般认为他是近代哲学的开端。
所以研究近代哲学，应当研究16世纪的历史变化。重要的变化有以下
几个：

\begin{itemize}
    \item 大航海时代（航路开辟，征服美洲）
    \item 文艺复兴，人文主义（文艺复兴高峰，北方文艺复兴）
    
    把握人文主义者的初衷和气质特征，反对中世纪发展成熟的技术化哲学
    \item 宗教改革（马丁·路德）
\end{itemize}

\subsection{对哲学发展的影响}
第16世纪的历史变化对哲学发展有深刻的影响：

\begin{itemize}
    \item 逻辑学的衰落
    \begin{enumerate}
        \item 模态逻辑研究被中断和遗忘
        \item 讲授简化版本的Aristotle逻辑
    \end{enumerate}
    \item 怀疑论再度流行
    \begin{enumerate}
        \item 塞尔维特案和卡斯特利奥的宽容主张
        \item 恩皮理科作品再度被发现
        \item 蒙田的怀疑论
    \end{enumerate}
    \item 现代科学的兴起
    \begin{enumerate}
        \item 伽利略（1564-1642）诸多观察和实验发现《两大世界体系的对话》1632：Salviati\&Simplicius《关于两种新科学的对话和数学证明》1638, 动力学和材料学
        \item Bacon（1561-1626）《新工具》：摈弃偏见；基于系统考查的归纳（种族，巢穴，市场，剧院）Induction：形成科学的视角（归纳法研究科学）；倡导广泛的科学合作
    \end{enumerate}
\end{itemize}

\begin{fancybox}{蒙田的怀疑论}
\begin{itemize}
    \item 与恩皮理科类似，借用伊壁鸠鲁和斯多亚学派的两方的资源
    \item 伊壁鸠鲁：如果感觉不可靠，那么没有知识\\斯多亚学派：知识不能来自感觉，因为感觉不可靠
    \item 一些常见的错觉例子
    \item 拔高人和其他物种
\end{itemize}
总结下来，所有感官是不可知的，因而所有知识是不可信的。
\end{fancybox}

接下来本章的思考题（本部分来源于课件）：

\begin{itemize}
    \item 影响哲学研究动态的因素有哪些？
    \item 一个人可以成为真诚和确定的怀疑论者吗？为什么
    \item 如果只能选取一个人作为现代科学的先驱，应该是Garleio还是Bacon？
\end{itemize}

